\documentclass[12pt]{article}
\usepackage{enumitem}
\usepackage{hyperref}
\usepackage{amsmath}
\usepackage[T1]{fontenc}

\title{Meu Relatório \LaTeX\ }
\author{Pedro H. Pereira da Silva}
\date{\today}

\begin{document}
\maketitle

\section{Introdução}
Neste relatório veremos resumidamente 3 conceitos vistos em aula no formato de lista. Além disso, iremos ver a aplicação da ferramenta \textbf{BibTeX}  que serve como gerenciador de referências bibliogŕaficas.

    \subsection{Conceitos}
    Seguimos então para os conceitos:
    \begin{enumerate}[label=\bf{\Roman*}.]
        \item \textbf{Shell:} \\
        Programa de sistemas Unix/Linux que forma interface entre usuários e o sistema. \\
        Por exemplo, o comando a seguir utiliza de \textit{pipes} (|) para buscar arquivos contendo os termos "prmpt" e "pdf" de forma \textit{case insensitive} dentre todos os arquivos listados na pasta com os termos encontrados coloridos.
        
        \texttt{ls | sort | grep --color=always -i prmpt | grep --color=always -i PDF | more}

        \item \textbf{Git:} \\
        É um sistema de controle de versão open-source. Através dessa ferramenta podemos controlar e organizar versões de códigos, além de poder fazer conexão direta de um repositório online com pasta local, alterando ambos através dos então chamados "commits".

        \item \textbf{LaTeX:} \\
        A ferramenta utilizada para escrever artigos científicos e a que estamos utilizando, o famoso \LaTeX. Através dele podemos utilizar diversos "truques" que diminuem imensamente a quantidade de horas de trabalho e de dor de cabeça. \\
        Vejamos algumas de suas utilidades:
        \begin{enumerate}[label=\roman*.]
            \item Hyperlinks: \\
            Através do pacote "\textit{hyperref}" podemos criar hyperlinks como: \href{https://github.com/dune-mdias/F056/tree/main}{Site da Matéria.}
            \item Expressões Matemáticas: \\
            \begin{equation}
            V_{eff} \ = \ -\frac{4}{3}\frac{\alpha_{s}(r)\hbar c}{r} \ + \ kr \\
            \end{equation}
            \begin{equation}
            \frac{\partial\rho}{\partial t} \ + \ \nabla \cdot \vec{J} \ =  \ 0 \\
            \end{equation}
            \begin{equation}
            i\hbar \frac{\partial}{\partial t} \psi(x,t) = \hat{H} \psi(x,t) \\
            \end{equation} 
            \begin{equation}
            \frac{dN}{d\phi} \propto 1 \ + \ 2\sum_{n=1}^{\infty}\nu_ncos[n(\phi-\Psi_n)] \
            \cite{KALWEIT:2023ll}
            \end{equation}
            \item Listas: \\
            Exatamente o que está sendo feito.
        \end{enumerate}
    \end{enumerate}

\section{Conclusão}
Agora faremos a conclusão. Para isso iremos trabalhar com troubleshooting de um problema que nós iremos criar.

\bibliographystyle{unsrt}
\bibliography{referencias}
\end{document}
